\documentclass[10pt,a4paper]{article}
\usepackage[margin=1.9cm]{geometry}
\usepackage[T1]{fontenc}
\usepackage{lmodern}
\usepackage{amsmath}
\usepackage{xcolor}
\usepackage{booktabs}
\usepackage{array}
\usepackage{enumitem}
\usepackage{titlesec}

\definecolor{fpgablue}{RGB}{13,71,161}
\definecolor{fpgalight}{RGB}{187,222,251}
\definecolor{hostgreen}{RGB}{27,94,32}
\definecolor{accent}{RGB}{183,28,28}
\definecolor{rulegray}{RGB}{120,120,120}
\definecolor{noteyellow}{RGB}{255,249,196}

\usepackage[colorlinks=true]{hyperref}
\hypersetup{linkcolor=fpgablue, urlcolor=blue,
            pdftitle={The Trade-Decision Brain}}

\setlength{\parindent}{0pt}
\setlength{\parskip}{5pt}

\titleformat{\section}{\large\bfseries\color{fpgablue}}{\thesection.}{0.5em}{}
\titleformat{\subsection}{\normalsize\bfseries\color{fpgablue}}{}{0em}{}
\titlespacing*{\section}{0pt}{10pt}{4pt}
\titlespacing*{\subsection}{0pt}{6pt}{2pt}

\newcommand{\hruleline}{\textcolor{rulegray}{\rule{\linewidth}{1pt}}}
\newcommand{\code}[1]{\texttt{\small #1}}

\begin{document}

% ============================ HEADER =======================================
{\LARGE\bfseries\color{fpgablue} The Trade-Decision Brain}\hfill
{\small\color{rulegray} HFT Project \textbullet\ July 2026}\\[-1pt]
{\large What the system checks \emph{before} it places a trade --- modelled on how
real HFT desks decide.}

\hruleline

\section{The big idea: a signal is not a trade}

The FPGA (and its software twin, \code{strategy\_sw.py}) produces a \textbf{signal}
from order-book imbalance: ``there are far more buyers than sellers right now,
so lean long.'' That signal is \emph{necessary} but not \emph{sufficient}. A real
trading firm never fires an order on a raw signal alone --- it runs the signal
through a sequence of sanity checks, any one of which can \textbf{veto} the trade
or \textbf{shrink} its size.

This document explains that sequence --- the ``brain'' that sits between the raw
signal and an actual order. Two design rules shaped it:

\begin{itemize}[leftmargin=1.4em, itemsep=1pt]
  \item \textbf{The fast core is untouched.} The imbalance signal in
  \code{strategy\_sw.py} stays bit-for-bit identical to the FPGA (that equivalence
  is guarded by a test). All the logic here is a \emph{software overlay} on top ---
  exactly how real systems separate a fast signal engine from slower risk/execution
  checks.
  \item \textbf{Every rejection is counted.} The system logs \emph{why} each signal
  was dropped, so you can watch the flow thin out from many signals down to few
  orders (the ``funnel'' in \S\ref{sec:funnel}).
\end{itemize}

\section{The pre-trade gauntlet (in order)}

A signal must survive \emph{every} gate below, top to bottom, to become an order.
Code: \code{host/hft\_logic.py}.

\subsection{Gate 1 --- Fair value (microprice)}
\textbf{Question:} is the signal confirmed by the true fair price, not just the
midpoint? \\
\textbf{Plain version:} don't pay the sticker price --- work out what the thing is
actually worth right now. The naive ``mid'' is halfway between bid and ask, but if
there is far more size resting on the bid, the real fair value is a little
\emph{higher} than the mid. That size-weighted fair value is the \textbf{microprice}:
\[
  \text{microprice} \;=\; \frac{P_{\text{bid}}\cdot S_{\text{ask}}
    \;+\; P_{\text{ask}}\cdot S_{\text{bid}}}{S_{\text{bid}} + S_{\text{ask}}}
\]
We only \textbf{buy} if the microprice sits above the mid by at least a small margin
(\code{min\_edge\_ticks}), and only \textbf{sell} if it sits below by that margin.
If the microprice disagrees with the signal, the trade is vetoed.

\subsection{Gate 2 --- Edge versus cost}
\textbf{Question:} is the expected profit bigger than the cost of trading? \\
\textbf{Plain version:} don't drive across town to save a dollar on gas. Every trade
pays a fee (modelled as \code{fee\_bps}, basis points of the trade value). If the
expected edge (how far the microprice is from the mid) does not clear the fee, skip
the trade. This is what separates a \emph{profitable} signal from a merely
\emph{correct} one.

\subsection{Gate 3 --- Market-condition gates}
\textbf{Question:} is the market safe to trade into right now? Three quick checks:
\begin{itemize}[leftmargin=1.4em, itemsep=1pt]
  \item \textbf{Volatility} --- stand down when short-term price swings spike
  (\code{max\_vol\_bps}). In fast markets you get ``picked off'' by faster players.
  \item \textbf{Liquidity} --- require a minimum number of shares at the top of book
  (\code{min\_book\_size}). Thin books mean bad fills.
  \item \textbf{Staleness} --- require a fresh quote (\code{max\_quote\_age\_s}).
  Never act on an old price.
\end{itemize}

\subsection{Gate 4 --- Inventory skew}
\textbf{Question:} does this trade help or hurt my current position? \\
\textbf{Plain version:} lean toward flat. Trades that \emph{reduce} the position
always pass. Trades that \emph{add} to it are blocked once the position grows past a
fraction of the per-symbol cap (\code{inventory\_soft\_frac}). This stops the book
from piling up one-sided.

\subsection{Gate 5 --- Portfolio exposure \& the kill switch}
\textbf{Question:} across \emph{all} symbols, am I within my risk budget --- and
have I lost too much today? \\
This gate consults the shared \textbf{portfolio} (\S\ref{sec:portfolio}):
\begin{itemize}[leftmargin=1.4em, itemsep=1pt]
  \item \textbf{Gross-exposure cap} --- the order size is clamped to the dollar
  headroom left under the firm-wide gross-notional limit.
  \item \textbf{Kill switch} --- if the day's losses (or drawdown from the peak)
  breach the limit, \emph{all trading stops}, on every symbol. This is the single
  most important safety brake a real desk has.
\end{itemize}

\section{The funnel: many signals, few orders}
\label{sec:funnel}

Because every veto is counted, a run prints a funnel. A representative 7-symbol
(Magnificent~7) simulation:

\begin{center}
\begin{tabular}{lr}
\toprule
\textbf{Stage} & \textbf{Count} \\
\midrule
Raw signals evaluated            & 202 \\
\quad vetoed --- edge vs cost    & 139 \\
\quad vetoed --- inventory skew  & 9 \\
\quad vetoed --- liquidity       & 8 \\
\textbf{Orders passed}           & \textbf{46} \\
\bottomrule
\end{tabular}
\end{center}

Reading it: of 202 raw signals, most (139) were not worth their cost, a few were
blocked for risk/liquidity reasons, and 46 became orders. \emph{That thinning is the
whole point} --- it is the difference between a naive strategy that trades on every
twitch and a disciplined one that trades only when several independent conditions
agree.

\section{The portfolio: PnL, exposure, and the brake}
\label{sec:portfolio}

Code: \code{host/portfolio.py}. One portfolio object is shared by all symbols and
answers the questions a desk asks continuously:
\begin{itemize}[leftmargin=1.4em, itemsep=1pt]
  \item \textbf{Position} --- signed shares held per symbol (+ long, $-$ short).
  \item \textbf{Realized PnL} --- profit booked when a trade reduces a position,
  using average-cost accounting.
  \item \textbf{Unrealized PnL} --- open profit on the current position, marked to
  the latest price.
  \item \textbf{Net exposure} --- long minus short (directional risk).
  \item \textbf{Gross exposure} --- long plus short (total capital at work).
  \item \textbf{Kill switch} --- trips on a daily-loss or drawdown breach.
\end{itemize}

\section{An honest note on where the edge really is}

If you run the simulation you will often see equity drift slightly \emph{negative}.
That is not a bug --- it is the truth about a strategy that \emph{crosses the spread}
(a ``taker''): after fees, a naive imbalance signal is roughly break-even to slightly
losing. Real HFT profit comes from one of two places this project is honest about:
\begin{enumerate}[leftmargin=1.4em, itemsep=1pt]
  \item \textbf{Speed} --- being first to react. That is \emph{this project's actual
  edge}: the FPGA computes the decision in \textbf{205\,ns} versus microseconds in
  software. The decision \emph{logic} here is a faithful model; the \emph{advantage}
  is latency.
  \item \textbf{Market-making} --- posting two-sided quotes to \emph{earn} the spread
  rather than pay it. That is a different strategy than the aggressive taker modelled
  here, and a natural future extension.
\end{enumerate}

\section{Where it lives in the code}

\begin{center}
\begin{tabular}{ll}
\toprule
\textbf{File} & \textbf{Role} \\
\midrule
\code{rtl/strategy\_imbalance.sv}   & the raw signal, in hardware (FPGA) \\
\code{host/strategy\_sw.py}         & the raw signal, bit-identical in software \\
\code{host/hft\_logic.py}           & the pre-trade gauntlet (Gates 1--5) \\
\code{host/portfolio.py}            & positions, PnL, exposure, kill switch \\
\code{host/live\_feed.py}           & glues quotes $\to$ signal $\to$ gauntlet $\to$ order \\
\code{host/tb\_hft\_logic.py}       & unit tests for the whole overlay \\
\bottomrule
\end{tabular}
\end{center}

\textbf{Try it:} \code{python host/live\_feed.py --universe mag7 --simulate}
(dry-run, no real orders --- prints the funnel and portfolio at the end).

\end{document}
